\section{中值定理：连接局部与整体的桥梁}
\label{sec:02-Calculus/03-Differentiation/04}

从北京开到天津，里程表 120 公里，用时 1.5 小时，平均时速 80 公里/小时。路上并不匀速：出城堵在二三十公里的时速上，高速段一度跑到 110，进城前又慢下来。

问一个看似简单的问题：这一路上，速度表的指针有没有至少一次正好指在 80？

直觉告诉我们``应该有''。但支撑这个直觉的往往是``连续量必然经过中间值''（介值定理）。然而，我们手上真正拥有的核心工具不是速度的连续性，而是\textbf{导数}——每个时刻的瞬时变化率。

导数的麻烦在于它只活在一个点上：\(f'(3)=5\) 对第 5 秒或第 10 秒发生的事情无法提供任何信息。一堆孤立点上的局部信息，凭什么推出跨越整段区间的整体结论？

中值定理就是回答这个问题的核心工具。它在``局部导数''与``整体平均变化率''之间建立了一座双向通道。

\subsection{基石：Fermat 引理与极值点特征}

要建立这座通道，首先需要解决一个更基础的问题：\textbf{函数在内部取得极值时，局部导数呈现什么特征？}

\begin{lemma}[Fermat 引理]
设函数 \(f\) 在点 \(c\) 的某个邻域内有定义，且在 \(c\) 处可导。若 \(f\) 在 \(c\) 处取得局部极值，则 \(f'(c) = 0\)。
\end{lemma}

\begin{proof}
不妨设 \(f\) 在 \(c\) 处取得局部最大值。根据定义，在 \(c\) 的充分小邻域内，恒有 \(f(x) \le f(c)\)，即 \(f(x) - f(c) \le 0\)。
\begin{enumerate}
    \item 当 \(x > c\) 时，自变量增量 \(x - c > 0\)，故差商
    \[
    \frac{f(x) - f(c)}{x - c} \le 0 \implies f'(c) = \lim_{x \to c^+} \frac{f(x) - f(c)}{x - c} \le 0.
    \]
    \item 当 \(x < c\) 时，自变量增量 \(x - c < 0\)，故差商
    \[
    \frac{f(x) - f(c)}{x - c} \ge 0 \implies f'(c) = \lim_{x \to c^-} \frac{f(x) - f(c)}{x - c} \ge 0.
    \]
\end{enumerate}
由于 \(f\) 在 \(c\) 处可导，两侧极限必然存在且相等，因此只能 \(f'(c) = 0\)。
\end{proof}

Fermat 引理是连接几何极值与代数导数的关键桥梁。它指出：\textbf{可微函数在内部极值点处的切线必然水平}。这不是因为``图像看着是平的''，而是因为一个数不可能既 \(\ge 0\) 又 \(\le 0\) 却不等于零。

\subsection{先看两端齐平的情形：Rolle 定理}

有了 Fermat 引理，我们先讨论最简单的情形：起点和终点高度一致（如绕操场跑一圈回到原点，净位移为零）。

\begin{theorem}[Rolle 定理]
若 \(f\) 在闭区间 \([a,b]\) 上连续，在开区间 \((a,b)\) 内可导，且两端点值相等 \(f(a)=f(b)\)，则存在至少一点 \(c\in(a,b)\)，使得 \(f'(c)=0\)。
\end{theorem}

Rolle 定理的证明逻辑非常清晰地展现了``整体保底，局部落脚''的思想：
\begin{enumerate}
    \item \textbf{整体存在性（闭区间连续性）}：由最值定理，连续函数在 \([a,b]\) 上必能取得最大值 \(M\) 与最小值 \(m\)。
    \item \textbf{分类讨论与局部一阶条件}：
    \begin{itemize}
        \item 若 \(M = m\)，则 \(f(x)\) 在 \([a,b]\) 上为常数函数，内部任意点的导数均为 \(0\)。
        \item 若 \(M \neq m\)，因 \(f(a)=f(b)\)，最值 \(M\) 与 \(m\) 中至少有一个在开区间内部 \((a,b)\) 的某点 \(c\) 处取得。由 Fermat 引理，该内部极值点必满足 \(f'(c)=0\)。
    \end{itemize}
\end{enumerate}

\subsection{倾斜割线的平移：Lagrange 中值定理}

回到开车的例子，起点与终点位置不同，即 \(f(a) \neq f(b)\)。此时两端连接成的割线斜率为 \(\frac{f(b)-f(a)}{b-a}\)（即平均速度）。

如何寻找切线平行于割线的点？

直观上看，如果我们沿垂直于割线的方向平移割线，直到它与曲线相切，切点就是曲线距离割线``竖直距离最远''的点。
因此，曲线与割线之间的竖直距离函数为：
\[
g(x) = f(x) - \left[ f(a) + \frac{f(b)-f(a)}{b-a}(x-a) \right].
\]
其中方括号内为端点割线的方程。显然，\(g(x)\) 衡量了曲线偏离割线的程度，且满足端点齐平条件 \(g(a) = g(b) = 0\)。

由于减去一次线性函数不改变连续性与可导性，\(g(x)\) 在 \([a,b]\) 上完全满足 Rolle 定理的假设条件。故存在 \(c \in (a,b)\) 使得 \(g'(c) = 0\)。求导即得：
\[
f'(c) - \frac{f(b)-f(a)}{b-a} = 0 \implies f'(c) = \frac{f(b)-f(a)}{b-a}.
\]

由此我们正式得到微分学的核心定理之一：

\begin{theorem}[Lagrange 中值定理]
若函数 \(f\) 满足以下条件：
\begin{enumerate}
    \item 在闭区间 \([a,b]\) 上连续；
    \item 在开区间 \((a,b)\) 内可导；
\end{enumerate}
则在开区间 \((a,b)\) 内至少存在一点 \(c \in (a,b)\)，使得
\[
f'(c) = \frac{f(b)-f(a)}{b-a},
\]
或等价地写为 \(f(b) - f(a) = f'(c)(b - a)\)。
\end{theorem}

回到开车的例子，位置函数满足在时间区间 \([0, 1.5]\) 上连续且在 \((0, 1.5)\) 内可导，这意味着途中必有某一时刻 \(c \in (0, 1.5)\)，速度表的指针恰好指向 80 公里/小时。

\begin{center}
\begin{tikzpicture}[x=0.8cm,y=1cm,font=\small]
  \draw[black!45] (-0.3,0) -- (5.9,0);
  \draw[blue!70!black,thick] plot[domain=0:5.5,samples=90] (\x,{0.45*\x+0.9*sin(\x r)});
  \draw[gray!75,thick] (0,0) -- (5.5,1.84);
  \draw[black,thick] (0.7,1.323) -- (2.7,1.992);
  \draw[black,thick] (3.6,0.841) -- (5.5,1.476);
  \fill (0,0) circle (2pt);
  \fill (5.5,1.84) circle (2pt);
  \fill (1.699,1.657) circle (2pt);
  \fill (4.584,1.170) circle (2pt);
  \node[below,font=\footnotesize] at (1.699,1.50) {\(c_1\)};
  \node[below,font=\footnotesize] at (4.584,1.02) {\(c_2\)};
  \node[above,gray!80,font=\footnotesize] at (4.1,1.32) {割线};
  \node[below,font=\footnotesize] at (0,-0.15) {\(a\)};
  \node[below,font=\footnotesize] at (5.5,-0.15) {\(b\)};
  \begin{scope}[shift={(7.6,0)}]
    \draw[black!45] (-0.3,0) -- (2.9,0);
    \draw[blue!70!black,thick] (0,1.2) -- (1.3,0) -- (2.6,1.2);
    \draw[gray!75,thick] (0,1.2) -- (2.6,1.2);
    \fill (1.3,0) circle (2pt);
    \node[below,font=\footnotesize] at (1.3,-0.15) {尖角};
    \node[above,gray!80,font=\footnotesize] at (1.3,1.2) {割线水平};
  \end{scope}
\end{tikzpicture}

\emph{左：条件满足时，内部可能存在多个平行的切点（定理仅保证存在性，不保证唯一性）。右：\(f(x)=|x|\) 在 \([-1,1]\) 上两端齐平，但不可导的``尖角''使得两侧差商极限不一致，定理结论不再成立。}
\end{center}

\subsection{假设条件的边界：每条约束如何防止定理失效}

Lagrange 中值定理的条件精炼到了极致，删去或削弱任何一条，结论都会崩溃：
\begin{itemize}
    \item \textbf{开区间 \((a,b)\) 内可导}：排除内部尖角（如 \(f(x)=|x|\) 在 \([-1,1]\)）。注意：\textbf{端点处不要求可导}（例如 \(f(x)=\sqrt{x}\) 在 \([0,1]\) 上，在 \(x=0\) 处右导数为 \(+\infty\)，但该函数在 \([0,1]\) 上连续且在 \((0,1)\) 内可导，定理结论依然成立）。
    \item \textbf{闭区间 \([a,b]\) 上连续}：排除端点跃迁或内部不连续。若端点不连续（如定义在 \([0,1]\) 上的 \(f(x)=x\)（\(0 \le x < 1\)）且 \(f(1)=2\)），割线斜率为 2，但内部任意点的导数均为 1，永远无法匹配——函数将变化量集中在端点``跳跃''完成，内部导数未能捕捉这一过程。
    \item \textbf{区间长度有限}：确保割线斜率 \(\frac{f(b)-f(a)}{b-a}\) 有明确的代数意义。
\end{itemize}

值得注意的是，中值定理仅承诺 \(c\) 的存在，并不给出具体位置。这种``不确定性''反而构成了它的强大之处：在多数理论推导中，我们无需关心 \(c\) 究竟在哪，只需要利用它的存在性，将代数关系精确地建立起来。

\subsection{应用示例：导数界与函数的全局控制}

中值定理最常见的应用形式是将等式转化为不等式。若在区间上导数受界限制，即 \(|f'(x)| \le M\)，则对任意两点 \(x_1, x_2\)，存在 \(c\) 使得：
\[
|f(x_2) - f(x_1)| = |f'(c)| \cdot |x_2 - x_1| \le M |x_2 - x_1|.
\]
导数的局部上界转化为函数整体变化的速率上限。

\textbf{数值估计示例}：在没有计算器的情况下估算 \(\sqrt{4.1}\)。
考察函数 \(f(x) = \sqrt{x}\) 在 \([4, 4.1]\) 上的变化。其导数为 \(f'(x) = \frac{1}{2\sqrt{x}}\)。
在区间 \([4, 4.1]\) 上，\(|f'(x)| \le \frac{1}{2\sqrt{4}} = 0.25\)。
由中值定理：
\[
|\sqrt{4.1} - 2| \le 0.25 \times (4.1 - 4) = 0.025.
\]
这表明 \(\sqrt{4.1}\) 与 2 的偏差不会超过 0.025（实际值为 \(\approx 2.02492\)，偏差约为 0.02492），界限非常紧凑。

\begin{remark}[高阶视角：从导数界到优化算法]
将上述不等式推广至多维空间（将 \(|f'|\) 替换为梯度的范数 \(\|\nabla f\|\)），该性质即为 \textbf{Lipschitz 连续性}。若对梯度本身再次应用中值定理，可得到``梯度 \(L\)-Lipschitz''假设。这是现代凸优化与一阶机器学习算法（如梯度下降）的核心前提：它限定了目标函数曲率的变化速率，直接决定了梯度下降算法步长（学习率）不能超过 \(\frac{2}{L}\)，否则会导致迭代震荡发散。
\end{remark}

\subsection{应用示例：根的存在性与唯一性拆解}

中值定理另一个直接推论是函数的单调性判据：若在区间上处处 \(f'(x) > 0\)，对任意 \(x_1 < x_2\)，由 \(f(x_2)-f(x_1) = f'(c)(x_2-x_1)\) 可知右端严格大于零，从而函数严格递增。

我们可以清晰地看到不同定理在解决``根的个数''问题上的分工：
考察方程 \(f(x) = x^3 + x - 1 = 0\) 在 \([0,1]\) 上的实根个数。
\begin{enumerate}
    \item \textbf{存在性}：\(f(0) = -1 < 0\)，\(f(1) = 1 > 0\)。由于 \(f\) 连续，由\textbf{介值定理}，区间内至少存在一个实根。
    \item \textbf{唯一性}：计算导数 \(f'(x) = 3x^2 + 1 \ge 1 > 0\)。由\textbf{中值定理}，\(f\) 在 \([0,1]\) 上严格递增，图像至多与 \(x\) 轴相交一次。
\end{enumerate}

组合上述两个独立逻辑，``存在且唯一''（恰有一个实根）的结论才得以完备确立。直观图像上的``单调曲线穿过 \(x\) 轴''掩盖了这两个定理各自独立的逻辑支撑；只有将其拆开，才能理解当函数的连续性或可导性被破坏时，究竟是存在性还是唯一性先告失效。